\section{Signal Model}
\label{sec:eft_signal_models}
\par
Both of the analyses described in this chapter were performed on the \{$isoscalar,isovector$\} basis.

This was chosen over the \{$neutron,proton$\} basis as it allows for comparison to both LUX \cite{LUX_RUN4_EFT_2021} and XENON \cite{xenon100_eft_ref} results and in future, Panda-X \cite{pandax_2_eft_ref}.
In this chapter, only $isoscalar$ signals are presented.

\par
The assumption used in this chapter is that a single operator dominates the interaction and that there is no mixing or interference between operators.
As such only one operator was considered at a time, with its coupling to nucleons set to 1 and all other operator couplings to 0.
In this case, the differential recoil rate introduced in \autoref{sec:eft_theory} scales linearly with the operator coupling.
This is just \autoref{eq:final_eft_differential_cross_section}:
\begin{equation}
    \frac{dR}{dE_R} = \frac{c^{0^2}_i \rho_\chi}{32 \pi m^3_\chi m^2_N} \int_{v_{min}} \frac{f(\vec{v})}{v} F^{0,0}_{i,i} (v^2, q^2) d^3 v
    \label{eq:analysis_eft_operator_signal_equation}
\end{equation} 

\par
The differential recoil spectra were generated using \textit{DMFormFactor}, a \textit{Mathematica} package \cite{dmformfactor_ref}.
Other software solutions exist to evaluate \autoref{eq:analysis_eft_operator_signal_equation} such as \textit{WimpyDD} (written in python) \cite{wimpydd_ref} and \textit{DMFortFactor} (written in fortran) \cite{dmfortfactor_ref}, however none of these provide exactly the same result due to the choice of normalisation, form factors and folding in of xenon isotope abundances used within each codebase.
\textit{DMFormFactor} was selected for this analysis as it is the same as is used by the LUX \cite{LUX_RUN4_EFT_2021}, Panda-X \cite{pandax_2_eft_ref} and XENON \cite{xenon100_eft_ref} collaborations.
The limits produced by XENON100 and LUX used \textit{DMFormFactor} v6.0.
The recent publication by Panda-X used an unpublished version which uses different nuclear form factors\footnote{unpublished at time of writing, however the public release is expected in early 2023}.
In order to maintain consistency with previous single operator limit analyses, \textit{DMFormFactor} v6.0 was used in the analyses presented in this chapter.

\par
Between the full exposure study and SR1, the parameters for the Standard Halo Model (SHM) were altered.
As mentioned in \autoref{chap:detection_theory}, the parameters for the SHM used for generating signal models within direct dark matter experiments have been standardised since mid-2021 \cite{standard_halo_model_conventions_ref}.
However, the full exposure study began before that date and so used the parameters previously used within the LZ collaboration, in \cite{LZ_projected_sensitivity_paper_ref,LZ_TechnicalDesignReview_ref,LZ_Ibles_LZStats_Thesis_ref}.
These are also the same parameters used in both the Xenon100 analysis and LUX RUN-4 analysis.
The SR1 study used the recommended astronomic properties.
Both sets of parameters are shown in \autoref{tab:DMFormFactor_parameters}.

\par
The differential recoil spectra for each operator were evaluated for WIMP masses of [5, 7, 10, 12, 14, 21, 33, 50, 100, 200, 400, 1000, 4000, 10000] GeV/c$^{2}$ out to 1000 keV.
The recoil spectra for WIMP masses of 50 GeV/c$^2$ and 1000 GeV/c$^2$ are shown in \autoref{fig:HENR_recoil_spectra_m50} and \autoref{fig:HENR_recoil_spectra_m1000} using the full exposure values.
The dips in the shape of the recoil spectra are due to the Bessel functions introduced for evaluating the form factors.

\par
In \autoref{fig:HENR_recoil_spectra_projected_vs_sr1}, a comparison between the SHM full exposure and SR1 values is shown for 2 operators.
As the dark matter velocity is larger in the SR1 case, the dark matter flux is increased.
Therefore the expected rate of recoils is greater.
In terms of sensitivity, the lack of seeing a signal with the SR1 parameters will set a tighter constraint on the WIMP-nucleon scattering.
A more complete comparison between how uncertainties in the standard model parameters affect sensitivity to dark matter searches can be found in \cite{LZ_Ibles_LZStats_Thesis_ref,billyboxer_thesis_ref}.

\begin{table}[]
    \centering
    \begin{tabular}{c|c|c}
        Parameter         & Projected Values   & SR1 Values       \\ \hline
        $\nu_0$           & 220 km/s           & 238 km/s         \\ 
        $\nu_{esc}$       & 544 km/s           & 544 km/s         \\
        $\rho_{\chi}$     & 0.3 GeV/cm$^{3}$   & 0.3 GeV/cm$^{3}$ \\
        $|\nu_E|$         & 245 km/s           & 252 km/s  
    \end{tabular}
    \caption{Standard Halo Model parameters used for full exposure and SR1 sensitivity studies.}
    \label{tab:DMFormFactor_parameters}
\end{table}


\begin{figure}[!htbp]%
\centering
\begin{tikzpicture}
\centering
  \begin{groupplot}[view={0}{90},
    group style = {group size = 1 by 3}]
    \iffalse
    \nextgroupplot[
    xlabel=Recoil Energy (keV),
    ylabel=Differential Rate (kg/day/keV),
    legend pos=north east,
    mark size=0pt,
    xmin=1, xmax=400, xmode=log,
    ymin=1e-14, ymax=1e3, ymode=log]
    \addplot table {Data/HENR/Projected_Sensitivity/Signal/Recoils/dRdEr_EFT_O01s_m50GeV_0.dat};
    \addplot table {Data/HENR/Projected_Sensitivity/Signal/Recoils/dRdEr_EFT_O11s_m50GeV_0.dat};       
    \legend{$\Operator$1s, $\Operator$11s};
    \fi
    \nextgroupplot[
    xlabel=Recoil Energy (keV),
    ylabel=Differential Rate (kg/day/keV),
    legend pos=north east,
    mark size=0pt,
    xmin=1, xmax=400, xmode=log,
    ymin=1e-14, ymax=1e3, ymode=log,
    cycle list name=color list]
    \addplot table {Data/HENR/Projected_Sensitivity/Signal/Recoils/dRdEr_EFT_O04s_m50GeV_0.dat};
    \addplot table {Data/HENR/Projected_Sensitivity/Signal/Recoils/dRdEr_EFT_O06s_m50GeV_0.dat};     
    \addplot table {Data/HENR/Projected_Sensitivity/Signal/Recoils/dRdEr_EFT_O07s_m50GeV_0.dat};
    \addplot table {Data/HENR/Projected_Sensitivity/Signal/Recoils/dRdEr_EFT_O09s_m50GeV_0.dat};
    \addplot table {Data/HENR/Projected_Sensitivity/Signal/Recoils/dRdEr_EFT_O10s_m50GeV_0.dat};
    \addplot table {Data/HENR/Projected_Sensitivity/Signal/Recoils/dRdEr_EFT_O14s_m50GeV_0.dat};
    \legend{$\Operator$4s, $\Operator$6s, $\Operator$7s, $\Operator$9s, $\Operator$10s, $\Operator$14s};
    \iffalse
    \nextgroupplot[
    xlabel=Recoil Energy (keV),
    ylabel=Differential Rate (kg/day/keV),
    legend pos=north east,
    mark size=0pt,
    xmin=1, xmax=400, xmode=log,
    ymin=1e-14, ymax=1e3, ymode=log]
    \addplot table {Data/HENR/Projected_Sensitivity/Signal/Recoils/dRdEr_EFT_O03s_m50GeV_0.dat};
    \addplot table {Data/HENR/Projected_Sensitivity/Signal/Recoils/dRdEr_EFT_O05s_m50GeV_0.dat};
    \addplot table {Data/HENR/Projected_Sensitivity/Signal/Recoils/dRdEr_EFT_O08s_m50GeV_0.dat};
    \addplot table {Data/HENR/Projected_Sensitivity/Signal/Recoils/dRdEr_EFT_O12s_m50GeV_0.dat};
    \addplot table {Data/HENR/Projected_Sensitivity/Signal/Recoils/dRdEr_EFT_O13s_m50GeV_0.dat};
    \addplot table {Data/HENR/Projected_Sensitivity/Signal/Recoils/dRdEr_EFT_O15s_m50GeV_0.dat};       
    \legend{$\Operator$3s, $\Operator$5s, $\Operator$8s, $\Operator$12s, $\Operator$13s, $\Operator$15s};
    \fi
  \end{groupplot}
\end{tikzpicture}
\caption{Differential recoil spectra for the fourteen non-relativistic EFT WIMP-nucleon operators with a WIMP mass of 50 GeV/c$^2$.
         \textbf{Top:} spin-independent \textbf{Middle:} spin-dependent \textbf{Bottom:} novel }
\label{fig:HENR_recoil_spectra_m50}
\end{figure}


\begin{figure}[!htbp]%
\centering
\begin{tikzpicture}
\centering
  \begin{groupplot}[view={0}{90},
    group style = {group size = 1 by 3}]
    
    \nextgroupplot[
    xlabel=Recoil Energy (keV),
    ylabel=Differential Rate (kg/day/keV),
    legend pos=north east,
    mark size=0pt,
    xmin=1, xmax=400, xmode=log,
    ymin=1e-14, ymax=1e3, ymode=log]
    \addplot table {Data/HENR/Projected_Sensitivity/Signal/Recoils/dRdEr_EFT_O01s_m33GeV_0.dat};
    \addplot table {Data/HENR/Projected_Sensitivity/Signal/Recoils/dRdEr_EFT_O11s_m33GeV_0.dat};       
    \legend{$\Operator$1s, $\Operator$11s};
    \iffalse
    \nextgroupplot[
    xlabel=Recoil Energy (keV),
    ylabel=Differential Rate (kg/day/keV),
    legend pos=north east,
    mark size=0pt,
    xmin=1, xmax=400, xmode=log,
    ymin=1e-14, ymax=1e3, ymode=log,
    cycle list name=color list]
    \addplot table {Data/HENR/Projected_Sensitivity/Signal/Recoils/dRdEr_EFT_O04s_m1000GeV_0.dat};
    \addplot table {Data/HENR/Projected_Sensitivity/Signal/Recoils/dRdEr_EFT_O06s_m1000GeV_0.dat};     
    \addplot table {Data/HENR/Projected_Sensitivity/Signal/Recoils/dRdEr_EFT_O07s_m1000GeV_0.dat};
    \addplot table {Data/HENR/Projected_Sensitivity/Signal/Recoils/dRdEr_EFT_O09s_m1000GeV_0.dat};
    \addplot table {Data/HENR/Projected_Sensitivity/Signal/Recoils/dRdEr_EFT_O10s_m1000GeV_0.dat};
    \addplot table {Data/HENR/Projected_Sensitivity/Signal/Recoils/dRdEr_EFT_O14s_m1000GeV_0.dat};
    \legend{$\Operator$4s, $\Operator$6s, $\Operator$7s, $\Operator$9s, $\Operator$10s, $\Operator$14s};
    
    \nextgroupplot[
    xlabel=Recoil Energy (keV),
    ylabel=Differential Rate (kg/day/keV),
    legend pos=north east,
    mark size=0pt,
    xmin=1, xmax=400, xmode=log,
    ymin=1e-14, ymax=1e3, ymode=log]
    \addplot table {Data/HENR/Projected_Sensitivity/Signal/Recoils/dRdEr_EFT_O03s_m1000GeV_0.dat};
    \addplot table {Data/HENR/Projected_Sensitivity/Signal/Recoils/dRdEr_EFT_O05s_m1000GeV_0.dat};
    \addplot table {Data/HENR/Projected_Sensitivity/Signal/Recoils/dRdEr_EFT_O08s_m1000GeV_0.dat};
    \addplot table {Data/HENR/Projected_Sensitivity/Signal/Recoils/dRdEr_EFT_O12s_m1000GeV_0.dat};
    \addplot table {Data/HENR/Projected_Sensitivity/Signal/Recoils/dRdEr_EFT_O13s_m1000GeV_0.dat};
    \addplot table {Data/HENR/Projected_Sensitivity/Signal/Recoils/dRdEr_EFT_O15s_m1000GeV_0.dat};       
    \legend{$\Operator$3s, $\Operator$5s, $\Operator$8s, $\Operator$12s, $\Operator$13s, $\Operator$15s};
    \fi
  \end{groupplot}
\end{tikzpicture}
\caption{Differential recoil spectra for the fourteen non-relativistic EFT WIMP-nucleon operators with a WIMP mass of 1000 GeV/c$^2$.
         \textbf{Top:} spin-independent \textbf{Middle:} spin-dependent \textbf{Bottom:} novel }
\label{fig:HENR_recoil_spectra_m1000}
\end{figure}

\iffalse

\begin{figure}[!htbp]%
\centering
\begin{tikzpicture}
\centering
  \begin{groupplot}[view={0}{90},
    group style = {group size = 3 by 1}]
    \nextgroupplot[
    xlabel=Recoil Energy (keV),
    ylabel=Differential Rate (kg/day/keV),
    legend pos=north east,
    mark size=0pt,
    xmin=1, xmax=400, xmode=log,
    ymin=1e-14, ymax=1e3, ymode=log]
    \foreach \henrop in \roguerecoiloperators{
                \addplot 
                    table
                    %{Data/HENR/Signal/Recoils/dRdEr_EFT_O01s_m5GeV_0.txt};
                    {Data/HENR/Projected_Sensitivity/Signal/Recoils/dRdEr_EFT_O\henrop_m5GeV_0.dat};
            }
    \nextgroupplot[
    xlabel=Recoil Energy (keV),
    legend pos=north east,
    mark size=0pt,
    xmin=1, xmax=400, xmode=log,
    ymin=1e-14, ymax=1e3, ymode=log]
    \foreach \henrop in \roguerecoiloperators{
                \addplot
                    table
                    %{Data/HENR/Signal/Recoils/dRdEr_EFT_O01s_m5GeV_0.txt};
                    {Data/HENR/Projected_Sensitivity/Signal/Recoils/dRdEr_EFT_O\henrop_m5GeV_0.dat};
            }
  \end{groupplot}
\end{tikzpicture}
\caption{Literal Aids
\textbf{Left:} $Gd^{156}$ de-excitation path.
\textbf{Right:} $Gd^{158}$ de-excitation path.
}
\label{fig:HENR_NotSpin_Recoil_Spectrum}
\end{figure}


\begin{figure}[!htbp]%
\centering
\begin{tikzpicture}
\centering
  \begin{groupplot}[view={0}{90},
    group style = {group size = 2 by 1}]
    \nextgroupplot[
    xlabel=Recoil Energy (keV),
    ylabel=Differential Rate (kg/day/keV),
    legend pos=north east,
    mark size=0pt,
    xmin=1, xmax=400, xmode=log,
    ymin=1e-14, ymax=1e3, ymode=log]
    \foreach \henrop in \spinrecoiloperators{
                \addplot 
                    table
                    %{Data/HENR/Signal/Recoils/dRdEr_EFT_O01s_m5GeV_0.txt};
                    {Data/HENR/Projected_Sensitivity/Signal/Recoils/dRdEr_EFT_O\henrop_m5GeV_0.dat};
            }
    \nextgroupplot[
    xlabel=Recoil Energy (keV),
    legend pos=north east,
    mark size=0pt,
    xmin=1, xmax=400, xmode=log,
    ymin=1e-14, ymax=1e3, ymode=log]
    \foreach \henrop in \spinrecoiloperators{
                \addplot
                    table
                    %{Data/HENR/Signal/Recoils/dRdEr_EFT_O01s_m5GeV_0.txt};
                    {Data/HENR/Projected_Sensitivity/Signal/Recoils/dRdEr_EFT_O\henrop_m21GeV_0.dat};
            }
  \end{groupplot}
\end{tikzpicture}
\caption{Literal Aids
\textbf{Left:} $Gd^{156}$ de-excitation path.
\textbf{Right:} $Gd^{158}$ de-excitation path.
}
\label{fig:HENR_Spin_Recoil_Spectrum}
\end{figure}

\begin{figure}[!htbp]%
\centering
\begin{tikzpicture}
\centering
  \begin{groupplot}[view={0}{90},
    group style = {group size = 2 by 1}]
    \nextgroupplot[
    xlabel=Recoil Energy (keV),
    ylabel=Differential Rate (kg/day/keV),
    legend pos=north east,
    mark size=0pt,
    xmin=1, xmax=400, xmode=log,
    ymin=1e-14, ymax=1e3, ymode=log]
    \foreach \henrop in \roguerecoiloperators{
                \addplot 
                    table
                    %{Data/HENR/Signal/Recoils/dRdEr_EFT_O01s_m5GeV_0.txt};
                    {Data/HENR/Projected_Sensitivity/Signal/Recoils/dRdEr_EFT_O\henrop_m5GeV_0.dat};
            }
    \nextgroupplot[
    xlabel=Recoil Energy (keV),
    legend pos=north east,
    mark size=0pt,
    xmin=1, xmax=400, xmode=log,
    ymin=1e-14, ymax=1e3, ymode=log]
    \foreach \henrop in \roguerecoiloperators{
                \addplot
                    table
                    %{Data/HENR/Signal/Recoils/dRdEr_EFT_O01s_m5GeV_0.txt};
                    {Data/HENR/Projected_Sensitivity/Signal/Recoils/dRdEr_EFT_O\henrop_m5GeV_0.dat};
            }
  \end{groupplot}
\end{tikzpicture}
\caption{Literal Aids
\textbf{Left:} $Gd^{156}$ de-excitation path.
\textbf{Right:} $Gd^{158}$ de-excitation path.
}
\label{fig:HENR_NotSpin_Recoil_Spectrum}
\end{figure}

\fi

\begin{figure} %[!htbp]%
\centering
\begin{tikzpicture}
\centering
  \begin{groupplot}[view={0}{90},
    group style = {group size = 1 by 2}]
    
    \nextgroupplot[
    %xlabel=Recoil Energy (keV),
    ylabel=Differential Rate (kg/day/keV),
    legend pos=north east,
    width=15cm,
    height=7cm,
    mark size=0pt,
    xmin=0, xmax=1000,
    ymin=1e-14, ymax=1e2, ymode=log,
    cycle list name=color list] 
    \addplot table [x=energy, y=rate] {Data/HENR/sr1_ws/Signal/recoils/recoil_operator01_m400.dat}; 
    \addplot table [x=energy, y=rate] {Data/HENR/Projected_Sensitivity/Signal/recoils/recoil_operator01_m400.dat};
    \legend{projected, sr1};
    
    \nextgroupplot[
    xlabel=Recoil Energy (keV),
    ylabel=Differential Rate (kg/day/keV),
    legend pos=north east,
    width=15cm,
    height=7cm,
    mark size=0pt,
    xmin=0, xmax=1000,
    ymin=1e-16, ymax=1e-6, ymode=log,
    cycle list name=color list]
    \addplot table [x=energy, y=rate] {Data/HENR/Projected_Sensitivity/Signal/recoils/recoil_operator15_m400.dat}; 
    \addplot table [x=energy, y=rate] {Data/HENR/sr1_ws/Signal/recoils/recoil_operator15_m400.dat}; 
    \legend{projected, sr1};
    
  \end{groupplot}
\end{tikzpicture}
\caption{Differential recoil spectra for isoscalar EFT operators for different standard halo model parameters.
         The WIMP mass is 400 GeV/c$^2$ in both plots.
         \textbf{Top:} $\Operator$1 \textbf{Bottom:} $\Operator$15 }
\label{fig:HENR_recoil_spectra_projected_vs_sr1}
\end{figure}


\iffalse
\par
Let us first discuss the signal model which we expect.
The differential rate we saw in CHAPTER XXX though difficult to calculate is fairly easy for the end user, with a suite of tools written to determine this.
Some of the original authors of the theory produced one such program, DMFormFactor (written in Mathematica) \cite{dmformfactor_ref}.
Other software solutions such as WimpyDD (python) and DMFortFactor (fortran) also exist.
None of these produce exactly the same rate which in large part is due to differing nuclear form factors and therefore different.

\par
The package used here is DMFormFactor v6.0 as was the choice for both LUX and XENON collaboration results\footnote{At the time of writing this is the newest public release. However these is a newer version which uses updated form factors which was used in \cite{pandax_2_eft_ref}.}.
The Standard Halo Model parameters used to describe the dark matter are summarised in \autoref{tab:DMFormFactor_parameters}.
\begin{table}[]
    \centering
    \begin{tabular}{c|c}
        Parameter         & Value  \\ \hline
        $\nu_0$           & 220$km s^{-1}$ \\
        $\nu_{esc}$       & 544$km s^{-1}$ \\
        $\rho_{\chi}$     & 0.3 $GeV/cm^{3}$ \\
        $|\nu_E|$         & 245 $km s^{-1}$ 
    \end{tabular}
    \caption{DMFormFactor parameters used for the standard halo model. CITE XXX}
    \label{tab:DMFormFactor_parameters}
\end{table}
To improve the realism of the recoil and to maintain consistency with previous experiments the natural abundance of Xe isotopes were folded in.
Each operator was generated out to 1000 keV and in the isoscalar basis and a selection of the resultant differential rates are shown in \autoref{fig:HENR_Spin_Recoil_Spectrum}.
Each Operator was performed for a WIMP mass of [5, 7, 10, 12, 14, 21, 33, 50, 100, 200, 400, 1000, 4000] GeV.

\begin{figure}[!htbp]%
\centering
\begin{tikzpicture}
\centering
  \begin{groupplot}[view={0}{90},
    group style = {group size = 1 by 3}]
    \iffalse
    \nextgroupplot[
    xlabel=Recoil Energy (keV),
    ylabel=Differential Rate (kg/day/keV),
    legend pos=north east,
    mark size=0pt,
    xmin=1, xmax=400, xmode=log,
    ymin=1e-14, ymax=1e3, ymode=log]
    \addplot table {Data/HENR/Projected_Sensitivity/Signal/Recoils/dRdEr_EFT_O01s_m50GeV_0.dat};
    \addplot table {Data/HENR/Projected_Sensitivity/Signal/Recoils/dRdEr_EFT_O11s_m50GeV_0.dat};       
    \legend{$\Operator$1s, $\Operator$11s};
    \fi
    \nextgroupplot[
    xlabel=Recoil Energy (keV),
    ylabel=Differential Rate (kg/day/keV),
    legend pos=north east,
    mark size=0pt,
    xmin=1, xmax=400, xmode=log,
    ymin=1e-14, ymax=1e3, ymode=log,
    cycle list name=color list]
    \addplot table {Data/HENR/Projected_Sensitivity/Signal/Recoils/dRdEr_EFT_O04s_m50GeV_0.dat};
    \addplot table {Data/HENR/Projected_Sensitivity/Signal/Recoils/dRdEr_EFT_O06s_m50GeV_0.dat};     
    \addplot table {Data/HENR/Projected_Sensitivity/Signal/Recoils/dRdEr_EFT_O07s_m50GeV_0.dat};
    \addplot table {Data/HENR/Projected_Sensitivity/Signal/Recoils/dRdEr_EFT_O09s_m50GeV_0.dat};
    \addplot table {Data/HENR/Projected_Sensitivity/Signal/Recoils/dRdEr_EFT_O10s_m50GeV_0.dat};
    \addplot table {Data/HENR/Projected_Sensitivity/Signal/Recoils/dRdEr_EFT_O14s_m50GeV_0.dat};
    \legend{$\Operator$4s, $\Operator$6s, $\Operator$7s, $\Operator$9s, $\Operator$10s, $\Operator$14s};
    \iffalse
    \nextgroupplot[
    xlabel=Recoil Energy (keV),
    ylabel=Differential Rate (kg/day/keV),
    legend pos=north east,
    mark size=0pt,
    xmin=1, xmax=400, xmode=log,
    ymin=1e-14, ymax=1e3, ymode=log]
    \addplot table {Data/HENR/Projected_Sensitivity/Signal/Recoils/dRdEr_EFT_O03s_m50GeV_0.dat};
    \addplot table {Data/HENR/Projected_Sensitivity/Signal/Recoils/dRdEr_EFT_O05s_m50GeV_0.dat};
    \addplot table {Data/HENR/Projected_Sensitivity/Signal/Recoils/dRdEr_EFT_O08s_m50GeV_0.dat};
    \addplot table {Data/HENR/Projected_Sensitivity/Signal/Recoils/dRdEr_EFT_O12s_m50GeV_0.dat};
    \addplot table {Data/HENR/Projected_Sensitivity/Signal/Recoils/dRdEr_EFT_O13s_m50GeV_0.dat};
    \addplot table {Data/HENR/Projected_Sensitivity/Signal/Recoils/dRdEr_EFT_O15s_m50GeV_0.dat};       
    \legend{$\Operator$3s, $\Operator$5s, $\Operator$8s, $\Operator$12s, $\Operator$13s, $\Operator$15s};
    \fi
  \end{groupplot}
\end{tikzpicture}
\caption{Differential recoil spectra for the fourteen non-relativistic EFT WIMP-nucleon operators with a WIMP mass of 50 GeV/c$^2$.
         \textbf{Top:} spin-independent \textbf{Middle:} spin-dependent \textbf{Bottom:} novel }
\label{fig:HENR_recoil_spectra_m50}
\end{figure}


\begin{figure}[!htbp]%
\centering
\begin{tikzpicture}
\centering
  \begin{groupplot}[view={0}{90},
    group style = {group size = 1 by 3}]
    
    \nextgroupplot[
    xlabel=Recoil Energy (keV),
    ylabel=Differential Rate (kg/day/keV),
    legend pos=north east,
    mark size=0pt,
    xmin=1, xmax=400, xmode=log,
    ymin=1e-14, ymax=1e3, ymode=log]
    \addplot table {Data/HENR/Projected_Sensitivity/Signal/Recoils/dRdEr_EFT_O01s_m33GeV_0.dat};
    \addplot table {Data/HENR/Projected_Sensitivity/Signal/Recoils/dRdEr_EFT_O11s_m33GeV_0.dat};       
    \legend{$\Operator$1s, $\Operator$11s};
    \iffalse
    \nextgroupplot[
    xlabel=Recoil Energy (keV),
    ylabel=Differential Rate (kg/day/keV),
    legend pos=north east,
    mark size=0pt,
    xmin=1, xmax=400, xmode=log,
    ymin=1e-14, ymax=1e3, ymode=log,
    cycle list name=color list]
    \addplot table {Data/HENR/Projected_Sensitivity/Signal/Recoils/dRdEr_EFT_O04s_m1000GeV_0.dat};
    \addplot table {Data/HENR/Projected_Sensitivity/Signal/Recoils/dRdEr_EFT_O06s_m1000GeV_0.dat};     
    \addplot table {Data/HENR/Projected_Sensitivity/Signal/Recoils/dRdEr_EFT_O07s_m1000GeV_0.dat};
    \addplot table {Data/HENR/Projected_Sensitivity/Signal/Recoils/dRdEr_EFT_O09s_m1000GeV_0.dat};
    \addplot table {Data/HENR/Projected_Sensitivity/Signal/Recoils/dRdEr_EFT_O10s_m1000GeV_0.dat};
    \addplot table {Data/HENR/Projected_Sensitivity/Signal/Recoils/dRdEr_EFT_O14s_m1000GeV_0.dat};
    \legend{$\Operator$4s, $\Operator$6s, $\Operator$7s, $\Operator$9s, $\Operator$10s, $\Operator$14s};
    
    \nextgroupplot[
    xlabel=Recoil Energy (keV),
    ylabel=Differential Rate (kg/day/keV),
    legend pos=north east,
    mark size=0pt,
    xmin=1, xmax=400, xmode=log,
    ymin=1e-14, ymax=1e3, ymode=log]
    \addplot table {Data/HENR/Projected_Sensitivity/Signal/Recoils/dRdEr_EFT_O03s_m1000GeV_0.dat};
    \addplot table {Data/HENR/Projected_Sensitivity/Signal/Recoils/dRdEr_EFT_O05s_m1000GeV_0.dat};
    \addplot table {Data/HENR/Projected_Sensitivity/Signal/Recoils/dRdEr_EFT_O08s_m1000GeV_0.dat};
    \addplot table {Data/HENR/Projected_Sensitivity/Signal/Recoils/dRdEr_EFT_O12s_m1000GeV_0.dat};
    \addplot table {Data/HENR/Projected_Sensitivity/Signal/Recoils/dRdEr_EFT_O13s_m1000GeV_0.dat};
    \addplot table {Data/HENR/Projected_Sensitivity/Signal/Recoils/dRdEr_EFT_O15s_m1000GeV_0.dat};       
    \legend{$\Operator$3s, $\Operator$5s, $\Operator$8s, $\Operator$12s, $\Operator$13s, $\Operator$15s};
    \fi
  \end{groupplot}
\end{tikzpicture}
\caption{Differential recoil spectra for the fourteen non-relativistic EFT WIMP-nucleon operators with a WIMP mass of 1000 GeV/c$^2$.
         \textbf{Top:} spin-independent \textbf{Middle:} spin-dependent \textbf{Bottom:} novel }
\label{fig:HENR_recoil_spectra_m1000}
\end{figure}

\iffalse

\begin{figure}[!htbp]%
\centering
\begin{tikzpicture}
\centering
  \begin{groupplot}[view={0}{90},
    group style = {group size = 3 by 1}]
    \nextgroupplot[
    xlabel=Recoil Energy (keV),
    ylabel=Differential Rate (kg/day/keV),
    legend pos=north east,
    mark size=0pt,
    xmin=1, xmax=400, xmode=log,
    ymin=1e-14, ymax=1e3, ymode=log]
    \foreach \henrop in \roguerecoiloperators{
                \addplot 
                    table
                    %{Data/HENR/Signal/Recoils/dRdEr_EFT_O01s_m5GeV_0.txt};
                    {Data/HENR/Projected_Sensitivity/Signal/Recoils/dRdEr_EFT_O\henrop_m5GeV_0.dat};
            }
    \nextgroupplot[
    xlabel=Recoil Energy (keV),
    legend pos=north east,
    mark size=0pt,
    xmin=1, xmax=400, xmode=log,
    ymin=1e-14, ymax=1e3, ymode=log]
    \foreach \henrop in \roguerecoiloperators{
                \addplot
                    table
                    %{Data/HENR/Signal/Recoils/dRdEr_EFT_O01s_m5GeV_0.txt};
                    {Data/HENR/Projected_Sensitivity/Signal/Recoils/dRdEr_EFT_O\henrop_m5GeV_0.dat};
            }
  \end{groupplot}
\end{tikzpicture}
\caption{Literal Aids
\textbf{Left:} $Gd^{156}$ de-excitation path.
\textbf{Right:} $Gd^{158}$ de-excitation path.
}
\label{fig:HENR_NotSpin_Recoil_Spectrum}
\end{figure}


\begin{figure}[!htbp]%
\centering
\begin{tikzpicture}
\centering
  \begin{groupplot}[view={0}{90},
    group style = {group size = 2 by 1}]
    \nextgroupplot[
    xlabel=Recoil Energy (keV),
    ylabel=Differential Rate (kg/day/keV),
    legend pos=north east,
    mark size=0pt,
    xmin=1, xmax=400, xmode=log,
    ymin=1e-14, ymax=1e3, ymode=log]
    \foreach \henrop in \spinrecoiloperators{
                \addplot 
                    table
                    %{Data/HENR/Signal/Recoils/dRdEr_EFT_O01s_m5GeV_0.txt};
                    {Data/HENR/Projected_Sensitivity/Signal/Recoils/dRdEr_EFT_O\henrop_m5GeV_0.dat};
            }
    \nextgroupplot[
    xlabel=Recoil Energy (keV),
    legend pos=north east,
    mark size=0pt,
    xmin=1, xmax=400, xmode=log,
    ymin=1e-14, ymax=1e3, ymode=log]
    \foreach \henrop in \spinrecoiloperators{
                \addplot
                    table
                    %{Data/HENR/Signal/Recoils/dRdEr_EFT_O01s_m5GeV_0.txt};
                    {Data/HENR/Projected_Sensitivity/Signal/Recoils/dRdEr_EFT_O\henrop_m21GeV_0.dat};
            }
  \end{groupplot}
\end{tikzpicture}
\caption{Literal Aids
\textbf{Left:} $Gd^{156}$ de-excitation path.
\textbf{Right:} $Gd^{158}$ de-excitation path.
}
\label{fig:HENR_Spin_Recoil_Spectrum}
\end{figure}

\begin{figure}[!htbp]%
\centering
\begin{tikzpicture}
\centering
  \begin{groupplot}[view={0}{90},
    group style = {group size = 2 by 1}]
    \nextgroupplot[
    xlabel=Recoil Energy (keV),
    ylabel=Differential Rate (kg/day/keV),
    legend pos=north east,
    mark size=0pt,
    xmin=1, xmax=400, xmode=log,
    ymin=1e-14, ymax=1e3, ymode=log]
    \foreach \henrop in \roguerecoiloperators{
                \addplot 
                    table
                    %{Data/HENR/Signal/Recoils/dRdEr_EFT_O01s_m5GeV_0.txt};
                    {Data/HENR/Projected_Sensitivity/Signal/Recoils/dRdEr_EFT_O\henrop_m5GeV_0.dat};
            }
    \nextgroupplot[
    xlabel=Recoil Energy (keV),
    legend pos=north east,
    mark size=0pt,
    xmin=1, xmax=400, xmode=log,
    ymin=1e-14, ymax=1e3, ymode=log]
    \foreach \henrop in \roguerecoiloperators{
                \addplot
                    table
                    %{Data/HENR/Signal/Recoils/dRdEr_EFT_O01s_m5GeV_0.txt};
                    {Data/HENR/Projected_Sensitivity/Signal/Recoils/dRdEr_EFT_O\henrop_m5GeV_0.dat};
            }
  \end{groupplot}
\end{tikzpicture}
\caption{Literal Aids
\textbf{Left:} $Gd^{156}$ de-excitation path.
\textbf{Right:} $Gd^{158}$ de-excitation path.
}
\label{fig:HENR_NotSpin_Recoil_Spectrum}
\end{figure}

\fi







\par
The actual signal model is actually produced by taking a theoretical event rate spectrum (produced by a Mathematica package - DMFormFactor - developed by XXX) and applying the analysis acceptance and detector response.
Passing through LZLama or PdfMaker, this turns the above into event observable - S1 and LogS2).
\par
DMFormFactor takes performs the calculations described previously.
Each Xenon isotope is evaluated separately and is weighted by the Xenon abundance before being added together to produce an energy spectrum.
As mentioned previously, the energy spectrum is dependant upon the choice of coupling and the WIMP mass.


\par
A subset of the recoil spectra calculated are shown in \autoref{fig:HENR_Spin_Recoil_Spectrum} and \autoref{fig:HENR_NotSpin_Recoil_Spectrum}.
The masses for which the energy spectra were calculated were;
all of which can be seen in Annex XXX.
The reason for this is that the shape of the limit produced is fairly well understood, so only a subset of WIMP masses are needed.

\begin{figure}[!htbp]%
\centering
\begin{tikzpicture}
\centering
  \begin{groupplot}[view={0}{90},
    group style = {group size = 1 by 3}]
    \iffalse
    \nextgroupplot[
    xlabel=Recoil Energy (keV),
    ylabel=Differential Rate (kg/day/keV),
    legend pos=north east,
    mark size=0pt,
    xmin=1, xmax=400, xmode=log,
    ymin=1e-14, ymax=1e3, ymode=log]
    \addplot table {Data/HENR/Projected_Sensitivity/Signal/Recoils/dRdEr_EFT_O01s_m50GeV_0.dat};
    \addplot table {Data/HENR/Projected_Sensitivity/Signal/Recoils/dRdEr_EFT_O11s_m50GeV_0.dat};       
    \legend{$\Operator$1s, $\Operator$11s};
    \fi
    \nextgroupplot[
    xlabel=Recoil Energy (keV),
    ylabel=Differential Rate (kg/day/keV),
    legend pos=north east,
    mark size=0pt,
    xmin=1, xmax=400, xmode=log,
    ymin=1e-14, ymax=1e3, ymode=log,
    cycle list name=color list]
    \addplot table {Data/HENR/Projected_Sensitivity/Signal/Recoils/dRdEr_EFT_O04s_m50GeV_0.dat};
    \addplot table {Data/HENR/Projected_Sensitivity/Signal/Recoils/dRdEr_EFT_O06s_m50GeV_0.dat};     
    \addplot table {Data/HENR/Projected_Sensitivity/Signal/Recoils/dRdEr_EFT_O07s_m50GeV_0.dat};
    \addplot table {Data/HENR/Projected_Sensitivity/Signal/Recoils/dRdEr_EFT_O09s_m50GeV_0.dat};
    \addplot table {Data/HENR/Projected_Sensitivity/Signal/Recoils/dRdEr_EFT_O10s_m50GeV_0.dat};
    \addplot table {Data/HENR/Projected_Sensitivity/Signal/Recoils/dRdEr_EFT_O14s_m50GeV_0.dat};
    \legend{$\Operator$4s, $\Operator$6s, $\Operator$7s, $\Operator$9s, $\Operator$10s, $\Operator$14s};
    \iffalse
    \nextgroupplot[
    xlabel=Recoil Energy (keV),
    ylabel=Differential Rate (kg/day/keV),
    legend pos=north east,
    mark size=0pt,
    xmin=1, xmax=400, xmode=log,
    ymin=1e-14, ymax=1e3, ymode=log]
    \addplot table {Data/HENR/Projected_Sensitivity/Signal/Recoils/dRdEr_EFT_O03s_m50GeV_0.dat};
    \addplot table {Data/HENR/Projected_Sensitivity/Signal/Recoils/dRdEr_EFT_O05s_m50GeV_0.dat};
    \addplot table {Data/HENR/Projected_Sensitivity/Signal/Recoils/dRdEr_EFT_O08s_m50GeV_0.dat};
    \addplot table {Data/HENR/Projected_Sensitivity/Signal/Recoils/dRdEr_EFT_O12s_m50GeV_0.dat};
    \addplot table {Data/HENR/Projected_Sensitivity/Signal/Recoils/dRdEr_EFT_O13s_m50GeV_0.dat};
    \addplot table {Data/HENR/Projected_Sensitivity/Signal/Recoils/dRdEr_EFT_O15s_m50GeV_0.dat};       
    \legend{$\Operator$3s, $\Operator$5s, $\Operator$8s, $\Operator$12s, $\Operator$13s, $\Operator$15s};
    \fi
  \end{groupplot}
\end{tikzpicture}
\caption{Differential recoil spectra for the fourteen non-relativistic EFT WIMP-nucleon operators with a WIMP mass of 50 GeV/c$^2$.
         \textbf{Top:} spin-independent \textbf{Middle:} spin-dependent \textbf{Bottom:} novel }
\label{fig:HENR_recoil_spectra_m50}
\end{figure}


\begin{figure}[!htbp]%
\centering
\begin{tikzpicture}
\centering
  \begin{groupplot}[view={0}{90},
    group style = {group size = 1 by 3}]
    
    \nextgroupplot[
    xlabel=Recoil Energy (keV),
    ylabel=Differential Rate (kg/day/keV),
    legend pos=north east,
    mark size=0pt,
    xmin=1, xmax=400, xmode=log,
    ymin=1e-14, ymax=1e3, ymode=log]
    \addplot table {Data/HENR/Projected_Sensitivity/Signal/Recoils/dRdEr_EFT_O01s_m33GeV_0.dat};
    \addplot table {Data/HENR/Projected_Sensitivity/Signal/Recoils/dRdEr_EFT_O11s_m33GeV_0.dat};       
    \legend{$\Operator$1s, $\Operator$11s};
    \iffalse
    \nextgroupplot[
    xlabel=Recoil Energy (keV),
    ylabel=Differential Rate (kg/day/keV),
    legend pos=north east,
    mark size=0pt,
    xmin=1, xmax=400, xmode=log,
    ymin=1e-14, ymax=1e3, ymode=log,
    cycle list name=color list]
    \addplot table {Data/HENR/Projected_Sensitivity/Signal/Recoils/dRdEr_EFT_O04s_m1000GeV_0.dat};
    \addplot table {Data/HENR/Projected_Sensitivity/Signal/Recoils/dRdEr_EFT_O06s_m1000GeV_0.dat};     
    \addplot table {Data/HENR/Projected_Sensitivity/Signal/Recoils/dRdEr_EFT_O07s_m1000GeV_0.dat};
    \addplot table {Data/HENR/Projected_Sensitivity/Signal/Recoils/dRdEr_EFT_O09s_m1000GeV_0.dat};
    \addplot table {Data/HENR/Projected_Sensitivity/Signal/Recoils/dRdEr_EFT_O10s_m1000GeV_0.dat};
    \addplot table {Data/HENR/Projected_Sensitivity/Signal/Recoils/dRdEr_EFT_O14s_m1000GeV_0.dat};
    \legend{$\Operator$4s, $\Operator$6s, $\Operator$7s, $\Operator$9s, $\Operator$10s, $\Operator$14s};
    
    \nextgroupplot[
    xlabel=Recoil Energy (keV),
    ylabel=Differential Rate (kg/day/keV),
    legend pos=north east,
    mark size=0pt,
    xmin=1, xmax=400, xmode=log,
    ymin=1e-14, ymax=1e3, ymode=log]
    \addplot table {Data/HENR/Projected_Sensitivity/Signal/Recoils/dRdEr_EFT_O03s_m1000GeV_0.dat};
    \addplot table {Data/HENR/Projected_Sensitivity/Signal/Recoils/dRdEr_EFT_O05s_m1000GeV_0.dat};
    \addplot table {Data/HENR/Projected_Sensitivity/Signal/Recoils/dRdEr_EFT_O08s_m1000GeV_0.dat};
    \addplot table {Data/HENR/Projected_Sensitivity/Signal/Recoils/dRdEr_EFT_O12s_m1000GeV_0.dat};
    \addplot table {Data/HENR/Projected_Sensitivity/Signal/Recoils/dRdEr_EFT_O13s_m1000GeV_0.dat};
    \addplot table {Data/HENR/Projected_Sensitivity/Signal/Recoils/dRdEr_EFT_O15s_m1000GeV_0.dat};       
    \legend{$\Operator$3s, $\Operator$5s, $\Operator$8s, $\Operator$12s, $\Operator$13s, $\Operator$15s};
    \fi
  \end{groupplot}
\end{tikzpicture}
\caption{Differential recoil spectra for the fourteen non-relativistic EFT WIMP-nucleon operators with a WIMP mass of 1000 GeV/c$^2$.
         \textbf{Top:} spin-independent \textbf{Middle:} spin-dependent \textbf{Bottom:} novel }
\label{fig:HENR_recoil_spectra_m1000}
\end{figure}

\iffalse

\begin{figure}[!htbp]%
\centering
\begin{tikzpicture}
\centering
  \begin{groupplot}[view={0}{90},
    group style = {group size = 3 by 1}]
    \nextgroupplot[
    xlabel=Recoil Energy (keV),
    ylabel=Differential Rate (kg/day/keV),
    legend pos=north east,
    mark size=0pt,
    xmin=1, xmax=400, xmode=log,
    ymin=1e-14, ymax=1e3, ymode=log]
    \foreach \henrop in \roguerecoiloperators{
                \addplot 
                    table
                    %{Data/HENR/Signal/Recoils/dRdEr_EFT_O01s_m5GeV_0.txt};
                    {Data/HENR/Projected_Sensitivity/Signal/Recoils/dRdEr_EFT_O\henrop_m5GeV_0.dat};
            }
    \nextgroupplot[
    xlabel=Recoil Energy (keV),
    legend pos=north east,
    mark size=0pt,
    xmin=1, xmax=400, xmode=log,
    ymin=1e-14, ymax=1e3, ymode=log]
    \foreach \henrop in \roguerecoiloperators{
                \addplot
                    table
                    %{Data/HENR/Signal/Recoils/dRdEr_EFT_O01s_m5GeV_0.txt};
                    {Data/HENR/Projected_Sensitivity/Signal/Recoils/dRdEr_EFT_O\henrop_m5GeV_0.dat};
            }
  \end{groupplot}
\end{tikzpicture}
\caption{Literal Aids
\textbf{Left:} $Gd^{156}$ de-excitation path.
\textbf{Right:} $Gd^{158}$ de-excitation path.
}
\label{fig:HENR_NotSpin_Recoil_Spectrum}
\end{figure}


\begin{figure}[!htbp]%
\centering
\begin{tikzpicture}
\centering
  \begin{groupplot}[view={0}{90},
    group style = {group size = 2 by 1}]
    \nextgroupplot[
    xlabel=Recoil Energy (keV),
    ylabel=Differential Rate (kg/day/keV),
    legend pos=north east,
    mark size=0pt,
    xmin=1, xmax=400, xmode=log,
    ymin=1e-14, ymax=1e3, ymode=log]
    \foreach \henrop in \spinrecoiloperators{
                \addplot 
                    table
                    %{Data/HENR/Signal/Recoils/dRdEr_EFT_O01s_m5GeV_0.txt};
                    {Data/HENR/Projected_Sensitivity/Signal/Recoils/dRdEr_EFT_O\henrop_m5GeV_0.dat};
            }
    \nextgroupplot[
    xlabel=Recoil Energy (keV),
    legend pos=north east,
    mark size=0pt,
    xmin=1, xmax=400, xmode=log,
    ymin=1e-14, ymax=1e3, ymode=log]
    \foreach \henrop in \spinrecoiloperators{
                \addplot
                    table
                    %{Data/HENR/Signal/Recoils/dRdEr_EFT_O01s_m5GeV_0.txt};
                    {Data/HENR/Projected_Sensitivity/Signal/Recoils/dRdEr_EFT_O\henrop_m21GeV_0.dat};
            }
  \end{groupplot}
\end{tikzpicture}
\caption{Literal Aids
\textbf{Left:} $Gd^{156}$ de-excitation path.
\textbf{Right:} $Gd^{158}$ de-excitation path.
}
\label{fig:HENR_Spin_Recoil_Spectrum}
\end{figure}

\begin{figure}[!htbp]%
\centering
\begin{tikzpicture}
\centering
  \begin{groupplot}[view={0}{90},
    group style = {group size = 2 by 1}]
    \nextgroupplot[
    xlabel=Recoil Energy (keV),
    ylabel=Differential Rate (kg/day/keV),
    legend pos=north east,
    mark size=0pt,
    xmin=1, xmax=400, xmode=log,
    ymin=1e-14, ymax=1e3, ymode=log]
    \foreach \henrop in \roguerecoiloperators{
                \addplot 
                    table
                    %{Data/HENR/Signal/Recoils/dRdEr_EFT_O01s_m5GeV_0.txt};
                    {Data/HENR/Projected_Sensitivity/Signal/Recoils/dRdEr_EFT_O\henrop_m5GeV_0.dat};
            }
    \nextgroupplot[
    xlabel=Recoil Energy (keV),
    legend pos=north east,
    mark size=0pt,
    xmin=1, xmax=400, xmode=log,
    ymin=1e-14, ymax=1e3, ymode=log]
    \foreach \henrop in \roguerecoiloperators{
                \addplot
                    table
                    %{Data/HENR/Signal/Recoils/dRdEr_EFT_O01s_m5GeV_0.txt};
                    {Data/HENR/Projected_Sensitivity/Signal/Recoils/dRdEr_EFT_O\henrop_m5GeV_0.dat};
            }
  \end{groupplot}
\end{tikzpicture}
\caption{Literal Aids
\textbf{Left:} $Gd^{156}$ de-excitation path.
\textbf{Right:} $Gd^{158}$ de-excitation path.
}
\label{fig:HENR_NotSpin_Recoil_Spectrum}
\end{figure}

\fi

\par
When performing this analysis, an important choice is the coupling choice.


\paragraph{Copied from Billy}
\par
To allow for ease of comparison to previous limit setting on the inelastic36
WIMP-nucleon EFT operators by the XENON collaboration [? ], this analysis was conducted in the isoscalar basis.37
In the isoscalar basis, the charge densities of the nucleons are effectively averaged such that the interaction becomes38
indiscriminate to the type of nucleon involved, even though both the isoscalar and proton-neutron basis provide insight.39
For this WIMP-nucleon EFT, the UV scale governing the physics is far higher than the energies that are probed in the40
experiment. At these lower energies, the UV interactions have to be reduced to effective ones, which is done at the u,41
d and s quark level. The assumption is made that the couplings to each quark at these scales are roughly equivalent.42
Therefore the mass differences of these quarks are negligible at the high-energy scale of the underlying physics, and the43
interaction would be isoscalar. By performing this analysis in the isoscalar basis, it is possible to test this assumption’s44
validity. Additionally, by using either an isoscalar or isovector basis, the target’s nuclear state can be considered as45
isospin symmetric; a property of the strong force that can aid in simplifying the analysis.

\begin{figure}
    \centering
    \includegraphics[width=0.5\textwidth]{Figures/Placeholder.png}
    \caption{Integrated rate for each operator for a 1000GeV DM particle}
    \label{fig:operator_integrated_rate}
\end{figure}

\begin{table}[]
    \centering
    \begin{tabular}{c|c}
        Parameter   & Value  \\ \hline
        $g_{1}$     & 0.119 \\
        $g_{2}$     & 79.1  
    \end{tabular}
    \caption{Key detector parameters for the LXe-TPC parameters as used in \cite{LZ_projected_sensitivity_paper_ref}}
    \label{tab:projected_sensitivity_detector_parameters}
\end{table}

\fi